%     Bewertung der Modellleistung und Unterschiede zwischen CICEVSE2024 und dem eigenen Datensatz (z.B. Flooding-Attacken werden besser erkannt, Scans schwieriger).
The evaluation of the model demonstrates that the proposed approach performs reliably across both the CICEVSE2024 dataset and our own dataset. In particular, flooding attacks were consistently detected with very high accuracy, confirming that the selected features provide clear and distinguishable patterns for this type of malicious activity. The strong performance against flooding attacks indicates that the LSTM model in combination with decision-tree-based feature selection is capable of capturing temporal dependencies and amplifying signal peaks that are characteristic of such attacks.\\
However, the detection of scanning activities proved to be more challenging. The reason for this lies in the more subtle nature of scan-related anomalies, which often lack distinctive peaks in the observed features.\\

When comparing the results of this work with related studies, particularly the study presented in \cite{lstm_feature_selection}, a strong alignment can be observed. Both approaches demonstrate that combining decision-tree-based feature selection with an LSTM classifier leads to consistently high classification performance, with F1-scores exceeding 0.85. Furthermore, when the dataset exhibits strong peaks in the selected features, as is the case with the CICEVSE2024 dataset—F1-scores can surpass 0.95. This confirms that decision tree classifiers provide a powerful and efficient method for selecting features when deploying LSTM-based models.\\
This work also extends the discussion to embedded-hardware environments. A charging station setup similar to the one described in \cite{EV_Security} and implemented here could execute the detection algorithm efficiently without exceeding hardware constraints. The evaluation indicates that the optimal number of selected features lies between 11 and 20, balancing accuracy with computational efficiency. The detection algorithm, as well as the supporting monitoring tools such as Prometheus and perf, are not resource-intensive, which underlines the suitability of this approach for real-world embedded systems.\\

%     Abgleich deines Modells und der Feature-Selektion mit verwandten Studien.
As highlighted in both this work and related literature, flooding attacks can be detected with high reliability under comparable hardware configurations. However, the robustness of the detection results depends strongly on the underlying system setup and hardware event characteristics. If the hardware architecture differs significantly or kernel-level event traces deviate, model performance may degrade. This emphasizes the need for adaptive feature engineering and retraining when applying the model to new or heterogeneous environments.\\
A central limitation of this approach is the dependency on hardware performance events, which are inherently influenced by the system environment and running background processes. This makes cross-system transferability challenging, as results obtained on one hardware platform may not generalize without restrictions to another.\\
Another important limitation is the “black-box” nature of the LSTM model. While the model achieves high accuracy, its internal decision-making processes are not transparent. This lack of interpretability makes it difficult to provide precise explanations for why certain events are classified as attacks, particularly under different system conditions. Consequently, the robustness of the model may suffer when faced with unseen environments or unexpected event distributions.\\

% Einschränkungen:
Despite these limitations, the work presented here provides a strong foundation for the application of intrusion detection systems in electric vehicle charging stations. The results demonstrate that the proposed combination of LSTM and decision-tree-based feature selection can detect attacks with high accuracy, especially flooding attacks, while maintaining efficiency suitable for embedded hardware.